% --- Archon physics formalization source begin ---
% archon:physics
% archon:covers PhyXMiniProblems/problem_phyx_mini_0270.lean
% archon:source-report reports/phyx_mini/problem_phyx_mini_0270.source.json

\chapter{Physics problem phyx\_mini\_0270}
\label{ch:PhyXMiniProblems_problem_phyx_mini_0270}

\paragraph{Problem source.}
\#\# Physical scenario

In figure, a \$2.50\textbackslash{} kg\$ disk of diameter \$D = 42.0\textbackslash{} cm\$ is supported by a rod of length \$L = 76.0\textbackslash{} cm\$ and negligible mass that is pivoted at its end. With the torsion spring connected, the rod is vertical at equilibrium.

\#\# Question

What is the torsion constant of the spring if the period of oscillation has been decreased by \$0.500\textbackslash{} s\$?

\#\# Answer choices

- A: A: \$16.5\textbackslash{} N \textbackslash{}cdot m/rad\$

- B: B: \$17.5\textbackslash{} N \textbackslash{}cdot m/rad\$

- C: C: \$18.5\textbackslash{} N \textbackslash{}cdot m/rad\$

- D: D: \$19.5\textbackslash{} N \textbackslash{}cdot m/rad\$

\#\# Figure caption (auxiliary; use the image as primary evidence)

The image depicts a pendulum system. Here's a breakdown of its components and details:

1. **Structure**: 
   - A horizontal bar at the top from which the pendulum hangs.
   - A rod extending downward from the bar.

2. **Objects**:
   - At the bottom of the rod, there is a circular disk representing the pendulum bob.

3. **Measurements**:
   - The vertical distance from the bottom of the horizontal bar to the center of the disk is labeled as \textbackslash{}( L \textbackslash{}).
   - The diameter of the disk is labeled as \textbackslash{}( D \textbackslash{}).

4. **Additional Elements**:
   - A circular marking at the top suggests a pivot or suspension point for the pendulum, with two concentric circles.
   - Arrows are used to denote measurements for \textbackslash{}( L \textbackslash{}) and \textbackslash{}( D \textbackslash{}).

This setup is typical of a simple pendulum used in physics to demonstrate oscillatory motion.

\#\# Dataset metadata

Resonance and Harmonics; Physical Model Grounding Reasoning, Multi-Formula Reasoning

\paragraph{Recorded answer/context.}
C: C: \$18.5\textbackslash{} N \textbackslash{}cdot m/rad\$

\paragraph{Figure/image path.}
/root/proposal\_for\_physic/science-mango/phyx\_mini\_run/phyx\_data/test\_image/270.png

\paragraph{Formalization target.}
create a compiling Lean file with sorry bodies at `PhyXMiniProblems/problem\_phyx\_mini\_0270.lean`. The Lean declarations must preserve the physical quantities, dimensions or dimensional roles, figure labels, governing-law hypotheses, and final relation expressed by this problem.
Use Mathlib/Physlib names found through LeanExplore where available. If a physics API is missing, introduce faithful local abstractions rather than scalar placeholder aliases.

\begin{theorem}[Physics formalization target]
\label{thm:physics:phyx_mini_0270:target}
\lean{PhyXMiniProblems.ProblemPhyXMini0270.diskTorsionPendulumConstant_matches_recordedAnswerC}
\uses{def:physics:phyx-mini-0270:phyxminiproblems-problemphyxmini0270-disktorsionpendulumsetup, def:physics:phyx-mini-0270:phyxminiproblems-problemphyxmini0270-matchesprimarydiskpendulumfigure, def:physics:phyx-mini-0270:phyxminiproblems-problemphyxmini0270-matchesdisktorsionpendulumproblemdata, def:physics:phyx-mini-0270:phyxminiproblems-problemphyxmini0270-hasphysicaldisktorsionpendulumparameters, def:physics:phyx-mini-0270:phyxminiproblems-problemphyxmini0270-satisfiesdisktorsionpendulumlaws, def:physics:phyx-mini-0270:phyxminiproblems-problemphyxmini0270-recordedanswerchoice, def:physics:phyx-mini-0270:phyxminiproblems-problemphyxmini0270-matchesdisplayedtorsionconstant, def:physics:phyx-mini-0270:phyxminiproblems-problemphyxmini0270-isuniqueclosesttorsionconstantchoice}
The assigned autoformalize agent should translate this physics problem into Lean declarations in the covered file, with theorem and lemma proof bodies written as `by sorry`.
\end{theorem}
\begin{proof}
This is an autoformalization task, not a proof task. Produce statements that can later be proved by the physics prover without weakening the physical model.
\end{proof}

% --- Archon declaration topology begin ---
\section{Declaration topology}
\begin{definition}[Mass Quantity]
  \label{def:physics:phyx-mini-0270:phyxminiproblems-problemphyxmini0270-massquantity}
  \lean{PhyXMiniProblems.ProblemPhyXMini0270.MassQuantity}
  A nonnegative physical mass
\end{definition}
\begin{proof}
  This is the declared model definition of Mass Quantity.
\end{proof}
\begin{definition}[Length Quantity]
  \label{def:physics:phyx-mini-0270:phyxminiproblems-problemphyxmini0270-lengthquantity}
  \lean{PhyXMiniProblems.ProblemPhyXMini0270.LengthQuantity}
  A nonnegative physical length
\end{definition}
\begin{proof}
  This is the declared model definition of Length Quantity.
\end{proof}
\begin{definition}[Time Quantity]
  \label{def:physics:phyx-mini-0270:phyxminiproblems-problemphyxmini0270-timequantity}
  \lean{PhyXMiniProblems.ProblemPhyXMini0270.TimeQuantity}
  A nonnegative physical duration
\end{definition}
\begin{proof}
  This is the declared model definition of Time Quantity.
\end{proof}
\begin{definition}[Acceleration Magnitude Quantity]
  \label{def:physics:phyx-mini-0270:phyxminiproblems-problemphyxmini0270-accelerationmagnitudequantity}
  \lean{PhyXMiniProblems.ProblemPhyXMini0270.AccelerationMagnitudeQuantity}
  A nonnegative acceleration magnitude, with dimension L T⁻²
\end{definition}
\begin{proof}
  This is the declared model definition of Acceleration Magnitude Quantity.
\end{proof}
\begin{definition}[Moment Of Inertia Quantity]
  \label{def:physics:phyx-mini-0270:phyxminiproblems-problemphyxmini0270-momentofinertiaquantity}
  \lean{PhyXMiniProblems.ProblemPhyXMini0270.MomentOfInertiaQuantity}
  A moment of inertia about the axis normal to the disk, with dimension M L²
\end{definition}
\begin{proof}
  This is the declared model definition of Moment Of Inertia Quantity.
\end{proof}
\begin{definition}[Torque Quantity]
  \label{def:physics:phyx-mini-0270:phyxminiproblems-problemphyxmini0270-torquequantity}
  \lean{PhyXMiniProblems.ProblemPhyXMini0270.TorqueQuantity}
  A signed torque, with dimension M L² T⁻²
\end{definition}
\begin{proof}
  This is the declared model definition of Torque Quantity.
\end{proof}
\begin{definition}[Restoring Torque Coefficient Quantity]
  \label{def:physics:phyx-mini-0270:phyxminiproblems-problemphyxmini0270-restoringtorquecoefficientquantity}
  \lean{PhyXMiniProblems.ProblemPhyXMini0270.RestoringTorqueCoefficientQuantity}
  The coefficient multiplying angular displacement in a linearized restoring torque. Radians are dimensionless, so its dimension is M L² T⁻²
\end{definition}
\begin{proof}
  This is the declared model definition of Restoring Torque Coefficient Quantity.
\end{proof}
\begin{definition}[Torsion Constant Quantity]
  \label{def:physics:phyx-mini-0270:phyxminiproblems-problemphyxmini0270-torsionconstantquantity}
  \lean{PhyXMiniProblems.ProblemPhyXMini0270.TorsionConstantQuantity}
  The torsion constant of the pivot spring. Although it has the same dimension as a restoring-torque coefficient, it has a distinct physical role
\end{definition}
\begin{proof}
  This is the declared model definition of Torsion Constant Quantity.
\end{proof}
\begin{definition}[mass In Kilograms]
  \label{def:physics:phyx-mini-0270:phyxminiproblems-problemphyxmini0270-massinkilograms}
  \lean{PhyXMiniProblems.ProblemPhyXMini0270.massInKilograms}
  \uses{def:physics:phyx-mini-0270:phyxminiproblems-problemphyxmini0270-massquantity}
  Read a physical mass in kilograms
\end{definition}
\begin{proof}
  This is the declared model definition of mass In Kilograms.
\end{proof}
\begin{definition}[length In Meters]
  \label{def:physics:phyx-mini-0270:phyxminiproblems-problemphyxmini0270-lengthinmeters}
  \lean{PhyXMiniProblems.ProblemPhyXMini0270.lengthInMeters}
  \uses{def:physics:phyx-mini-0270:phyxminiproblems-problemphyxmini0270-lengthquantity}
  Read a physical length in meters
\end{definition}
\begin{proof}
  This is the declared model definition of length In Meters.
\end{proof}
\begin{definition}[time In Seconds]
  \label{def:physics:phyx-mini-0270:phyxminiproblems-problemphyxmini0270-timeinseconds}
  \lean{PhyXMiniProblems.ProblemPhyXMini0270.timeInSeconds}
  \uses{def:physics:phyx-mini-0270:phyxminiproblems-problemphyxmini0270-timequantity}
  Read a physical duration in seconds
\end{definition}
\begin{proof}
  This is the declared model definition of time In Seconds.
\end{proof}
\begin{definition}[acceleration In Meters Per Second Squared]
  \label{def:physics:phyx-mini-0270:phyxminiproblems-problemphyxmini0270-accelerationinmeterspersecondsquared}
  \lean{PhyXMiniProblems.ProblemPhyXMini0270.accelerationInMetersPerSecondSquared}
  \uses{def:physics:phyx-mini-0270:phyxminiproblems-problemphyxmini0270-accelerationmagnitudequantity}
  Read an acceleration magnitude in meters per second squared
\end{definition}
\begin{proof}
  This is the declared model definition of acceleration In Meters Per Second Squared.
\end{proof}
\begin{definition}[moment Of Inertia In Kilogram Meters Squared]
  \label{def:physics:phyx-mini-0270:phyxminiproblems-problemphyxmini0270-momentofinertiainkilogrammeterssquared}
  \lean{PhyXMiniProblems.ProblemPhyXMini0270.momentOfInertiaInKilogramMetersSquared}
  \uses{def:physics:phyx-mini-0270:phyxminiproblems-problemphyxmini0270-momentofinertiaquantity}
  Read a moment of inertia in kilogram meter squared
\end{definition}
\begin{proof}
  This is the declared model definition of moment Of Inertia In Kilogram Meters Squared.
\end{proof}
\begin{definition}[torque In Newton Meters]
  \label{def:physics:phyx-mini-0270:phyxminiproblems-problemphyxmini0270-torqueinnewtonmeters}
  \lean{PhyXMiniProblems.ProblemPhyXMini0270.torqueInNewtonMeters}
  \uses{def:physics:phyx-mini-0270:phyxminiproblems-problemphyxmini0270-torquequantity}
  Read a signed torque in newton meters
\end{definition}
\begin{proof}
  This is the declared model definition of torque In Newton Meters.
\end{proof}
\begin{definition}[restoring Coefficient In Newton Meters Per Radian]
  \label{def:physics:phyx-mini-0270:phyxminiproblems-problemphyxmini0270-restoringcoefficientinnewtonmetersperradian}
  \lean{PhyXMiniProblems.ProblemPhyXMini0270.restoringCoefficientInNewtonMetersPerRadian}
  \uses{def:physics:phyx-mini-0270:phyxminiproblems-problemphyxmini0270-restoringtorquecoefficientquantity}
  Read a linearized restoring coefficient in newton meters per radian
\end{definition}
\begin{proof}
  This is the declared model definition of restoring Coefficient In Newton Meters Per Radian.
\end{proof}
\begin{definition}[torsion Constant In Newton Meters Per Radian]
  \label{def:physics:phyx-mini-0270:phyxminiproblems-problemphyxmini0270-torsionconstantinnewtonmetersperradian}
  \lean{PhyXMiniProblems.ProblemPhyXMini0270.torsionConstantInNewtonMetersPerRadian}
  \uses{def:physics:phyx-mini-0270:phyxminiproblems-problemphyxmini0270-torsionconstantquantity}
  Read a torsion constant in newton meters per radian
\end{definition}
\begin{proof}
  This is the declared model definition of torsion Constant In Newton Meters Per Radian.
\end{proof}
\begin{definition}[Disk Mass Distribution]
  \label{def:physics:phyx-mini-0270:phyxminiproblems-problemphyxmini0270-diskmassdistribution}
  \lean{PhyXMiniProblems.ProblemPhyXMini0270.DiskMassDistribution}
  The idealized mass distribution of the disk
\end{definition}
\begin{proof}
  This is the declared model definition of Disk Mass Distribution.
\end{proof}
\begin{definition}[Rod Mass Model]
  \label{def:physics:phyx-mini-0270:phyxminiproblems-problemphyxmini0270-rodmassmodel}
  \lean{PhyXMiniProblems.ProblemPhyXMini0270.RodMassModel}
  The mass model specified for the supporting rod
\end{definition}
\begin{proof}
  This is the declared model definition of Rod Mass Model.
\end{proof}
\begin{definition}[Equilibrium Orientation]
  \label{def:physics:phyx-mini-0270:phyxminiproblems-problemphyxmini0270-equilibriumorientation}
  \lean{PhyXMiniProblems.ProblemPhyXMini0270.EquilibriumOrientation}
  The equilibrium orientation stated in the problem
\end{definition}
\begin{proof}
  This is the declared model definition of Equilibrium Orientation.
\end{proof}
\begin{definition}[Pivot Spring Kind]
  \label{def:physics:phyx-mini-0270:phyxminiproblems-problemphyxmini0270-pivotspringkind}
  \lean{PhyXMiniProblems.ProblemPhyXMini0270.PivotSpringKind}
  The connected restoring element visible around the pivot
\end{definition}
\begin{proof}
  This is the declared model definition of Pivot Spring Kind.
\end{proof}
\begin{definition}[Figure Component]
  \label{def:physics:phyx-mini-0270:phyxminiproblems-problemphyxmini0270-figurecomponent}
  \lean{PhyXMiniProblems.ProblemPhyXMini0270.FigureComponent}
  Components visible in the supplied primary image
\end{definition}
\begin{proof}
  This is the declared model definition of Figure Component.
\end{proof}
\begin{definition}[Figure Point]
  \label{def:physics:phyx-mini-0270:phyxminiproblems-problemphyxmini0270-figurepoint}
  \lean{PhyXMiniProblems.ProblemPhyXMini0270.FigurePoint}
  Distinguished points needed to state the primary-image geometry
\end{definition}
\begin{proof}
  This is the declared model definition of Figure Point.
\end{proof}
\begin{definition}[Figure Label]
  \label{def:physics:phyx-mini-0270:phyxminiproblems-problemphyxmini0270-figurelabel}
  \lean{PhyXMiniProblems.ProblemPhyXMini0270.FigureLabel}
  Mathematical labels printed in the supplied image
\end{definition}
\begin{proof}
  This is the declared model definition of Figure Label.
\end{proof}
\begin{definition}[Disk Pendulum Figure]
  \label{def:physics:phyx-mini-0270:phyxminiproblems-problemphyxmini0270-diskpendulumfigure}
  \lean{PhyXMiniProblems.ProblemPhyXMini0270.DiskPendulumFigure}
  \uses{def:physics:phyx-mini-0270:phyxminiproblems-problemphyxmini0270-lengthquantity, def:physics:phyx-mini-0270:phyxminiproblems-problemphyxmini0270-figurecomponent, def:physics:phyx-mini-0270:phyxminiproblems-problemphyxmini0270-figurepoint, def:physics:phyx-mini-0270:phyxminiproblems-problemphyxmini0270-figurelabel}
  Structured readout of the supplied physical diagram
\end{definition}
\begin{proof}
  This is the declared model definition of Disk Pendulum Figure.
\end{proof}
\begin{definition}[Disk Torsion Pendulum Setup]
  \label{def:physics:phyx-mini-0270:phyxminiproblems-problemphyxmini0270-disktorsionpendulumsetup}
  \lean{PhyXMiniProblems.ProblemPhyXMini0270.DiskTorsionPendulumSetup}
  \uses{def:physics:phyx-mini-0270:phyxminiproblems-problemphyxmini0270-massquantity, def:physics:phyx-mini-0270:phyxminiproblems-problemphyxmini0270-lengthquantity, def:physics:phyx-mini-0270:phyxminiproblems-problemphyxmini0270-timequantity, def:physics:phyx-mini-0270:phyxminiproblems-problemphyxmini0270-accelerationmagnitudequantity, def:physics:phyx-mini-0270:phyxminiproblems-problemphyxmini0270-momentofinertiaquantity, def:physics:phyx-mini-0270:phyxminiproblems-problemphyxmini0270-torquequantity, def:physics:phyx-mini-0270:phyxminiproblems-problemphyxmini0270-restoringtorquecoefficientquantity, def:physics:phyx-mini-0270:phyxminiproblems-problemphyxmini0270-torsionconstantquantity, def:physics:phyx-mini-0270:phyxminiproblems-problemphyxmini0270-diskmassdistribution, def:physics:phyx-mini-0270:phyxminiproblems-problemphyxmini0270-rodmassmodel, def:physics:phyx-mini-0270:phyxminiproblems-problemphyxmini0270-equilibriumorientation, def:physics:phyx-mini-0270:phyxminiproblems-problemphyxmini0270-pivotspringkind, def:physics:phyx-mini-0270:phyxminiproblems-problemphyxmini0270-diskpendulumfigure}
  The physical system and its two periods. The spring torsion constant is an independent field: it is not defined from a displayed answer or from any other measurement. The signed torque functions describe the nonlinear system as a function of angular displacement in radians. The two Physlib oscillators and their periods represent the tangent systems at the vertical equilibrium, rather than asserting that the nonlinear finite-amplitude periods are exactly those of harmonic oscillators
\end{definition}
\begin{proof}
  This is the declared model definition of Disk Torsion Pendulum Setup.
\end{proof}
\begin{definition}[Matches Primary Disk Pendulum Figure]
  \label{def:physics:phyx-mini-0270:phyxminiproblems-problemphyxmini0270-matchesprimarydiskpendulumfigure}
  \lean{PhyXMiniProblems.ProblemPhyXMini0270.MatchesPrimaryDiskPendulumFigure}
  \uses{def:physics:phyx-mini-0270:phyxminiproblems-problemphyxmini0270-lengthinmeters, def:physics:phyx-mini-0270:phyxminiproblems-problemphyxmini0270-disktorsionpendulumsetup}
  Primary-image evidence. The L bracket runs from the pivot level to the top rim, while the D bracket spans the disk. Hence the center lever arm is L + D/2; this is figure geometry, not the requested torsion constant
\end{definition}
\begin{proof}
  This is the declared model definition of Matches Primary Disk Pendulum Figure.
\end{proof}
\begin{definition}[Matches Disk Torsion Pendulum Problem Data]
  \label{def:physics:phyx-mini-0270:phyxminiproblems-problemphyxmini0270-matchesdisktorsionpendulumproblemdata}
  \lean{PhyXMiniProblems.ProblemPhyXMini0270.MatchesDiskTorsionPendulumProblemData}
  \uses{def:physics:phyx-mini-0270:phyxminiproblems-problemphyxmini0270-massinkilograms, def:physics:phyx-mini-0270:phyxminiproblems-problemphyxmini0270-lengthinmeters, def:physics:phyx-mini-0270:phyxminiproblems-problemphyxmini0270-timeinseconds, def:physics:phyx-mini-0270:phyxminiproblems-problemphyxmini0270-accelerationinmeterspersecondsquared, def:physics:phyx-mini-0270:phyxminiproblems-problemphyxmini0270-disktorsionpendulumsetup}
  Numerical and qualitative problem data. The gravity value is the conventional near-Earth textbook calibration needed to distinguish the numerical choices. The period-reduction relation is an observed comparison between two physical states, interpreted by the textbook problem as a comparison of the two small-angle tangent-model periods. It does not assign a value to the spring constant
\end{definition}
\begin{proof}
  This is the declared model definition of Matches Disk Torsion Pendulum Problem Data.
\end{proof}
\begin{definition}[Has Physical Disk Torsion Pendulum Parameters]
  \label{def:physics:phyx-mini-0270:phyxminiproblems-problemphyxmini0270-hasphysicaldisktorsionpendulumparameters}
  \lean{PhyXMiniProblems.ProblemPhyXMini0270.HasPhysicalDiskTorsionPendulumParameters}
  \uses{def:physics:phyx-mini-0270:phyxminiproblems-problemphyxmini0270-massinkilograms, def:physics:phyx-mini-0270:phyxminiproblems-problemphyxmini0270-lengthinmeters, def:physics:phyx-mini-0270:phyxminiproblems-problemphyxmini0270-timeinseconds, def:physics:phyx-mini-0270:phyxminiproblems-problemphyxmini0270-accelerationinmeterspersecondsquared, def:physics:phyx-mini-0270:phyxminiproblems-problemphyxmini0270-momentofinertiainkilogrammeterssquared, def:physics:phyx-mini-0270:phyxminiproblems-problemphyxmini0270-restoringcoefficientinnewtonmetersperradian, def:physics:phyx-mini-0270:phyxminiproblems-problemphyxmini0270-torsionconstantinnewtonmetersperradian, def:physics:phyx-mini-0270:phyxminiproblems-problemphyxmini0270-disktorsionpendulumsetup}
  Positivity and nondegeneracy conditions for the physical system
\end{definition}
\begin{proof}
  This is the declared model definition of Has Physical Disk Torsion Pendulum Parameters.
\end{proof}
\begin{definition}[Satisfies Disk Torsion Pendulum Laws]
  \label{def:physics:phyx-mini-0270:phyxminiproblems-problemphyxmini0270-satisfiesdisktorsionpendulumlaws}
  \lean{PhyXMiniProblems.ProblemPhyXMini0270.SatisfiesDiskTorsionPendulumLaws}
  \uses{def:physics:phyx-mini-0270:phyxminiproblems-problemphyxmini0270-massinkilograms, def:physics:phyx-mini-0270:phyxminiproblems-problemphyxmini0270-lengthinmeters, def:physics:phyx-mini-0270:phyxminiproblems-problemphyxmini0270-timeinseconds, def:physics:phyx-mini-0270:phyxminiproblems-problemphyxmini0270-accelerationinmeterspersecondsquared, def:physics:phyx-mini-0270:phyxminiproblems-problemphyxmini0270-momentofinertiainkilogrammeterssquared, def:physics:phyx-mini-0270:phyxminiproblems-problemphyxmini0270-torqueinnewtonmeters, def:physics:phyx-mini-0270:phyxminiproblems-problemphyxmini0270-restoringcoefficientinnewtonmetersperradian, def:physics:phyx-mini-0270:phyxminiproblems-problemphyxmini0270-torsionconstantinnewtonmetersperradian, def:physics:phyx-mini-0270:phyxminiproblems-problemphyxmini0270-disktorsionpendulumsetup}
  Governing geometry, rigid-body, and small-oscillation laws: the radius is half the diameter; a uniform solid disk has central inertia m r² / 2; the parallel-axis contribution is m d²; the exact gravitational torque is -m g d sin θ; the ideal torsion spring obeys the exact Hooke law -κ θ; HasDerivAt contracts identify the tangent restoring coefficients at the vertical equilibrium, making the small-angle step local rather than a global equality;
\end{definition}
\begin{proof}
  This is the declared model definition of Satisfies Disk Torsion Pendulum Laws.
\end{proof}
\begin{lemma}[spring Torsion Constant from connected Period]
  \label{lem:physics:phyx-mini-0270:phyxminiproblems-problemphyxmini0270-springtorsionconstant-from-connectedperiod}
  \lean{PhyXMiniProblems.ProblemPhyXMini0270.springTorsionConstant_from_connectedPeriod}
  \uses{def:physics:phyx-mini-0270:phyxminiproblems-problemphyxmini0270-timeinseconds, def:physics:phyx-mini-0270:phyxminiproblems-problemphyxmini0270-momentofinertiainkilogrammeterssquared, def:physics:phyx-mini-0270:phyxminiproblems-problemphyxmini0270-restoringcoefficientinnewtonmetersperradian, def:physics:phyx-mini-0270:phyxminiproblems-problemphyxmini0270-torsionconstantinnewtonmetersperradian, def:physics:phyx-mini-0270:phyxminiproblems-problemphyxmini0270-disktorsionpendulumsetup, def:physics:phyx-mini-0270:phyxminiproblems-problemphyxmini0270-hasphysicaldisktorsionpendulumparameters, def:physics:phyx-mini-0270:phyxminiproblems-problemphyxmini0270-satisfiesdisktorsionpendulumlaws}
  Solving the spring-connected generalized-oscillator period law for the torsion constant gives κ = I (2π/T\_with)² - m g d. This remains a derived conclusion rather than a field of the governing-law interface
\end{lemma}
\begin{proof}
  The conclusion follows from the governing relations and figure readouts named in the statement of spring Torsion Constant from connected Period.
\end{proof}
\begin{definition}[Answer Choice]
  \label{def:physics:phyx-mini-0270:phyxminiproblems-problemphyxmini0270-answerchoice}
  \lean{PhyXMiniProblems.ProblemPhyXMini0270.AnswerChoice}
  Labels of the four displayed answer choices
\end{definition}
\begin{proof}
  This is the declared model definition of Answer Choice.
\end{proof}
\begin{definition}[newton Meters Per Radian]
  \label{def:physics:phyx-mini-0270:phyxminiproblems-problemphyxmini0270-answerchoice-newtonmetersperradian}
  \lean{PhyXMiniProblems.ProblemPhyXMini0270.AnswerChoice.newtonMetersPerRadian}
  \uses{def:physics:phyx-mini-0270:phyxminiproblems-problemphyxmini0270-answerchoice}
  Displayed torsion constants, in newton meters per radian
\end{definition}
\begin{proof}
  This is the declared model definition of newton Meters Per Radian.
\end{proof}
\begin{definition}[recorded Answer Choice]
  \label{def:physics:phyx-mini-0270:phyxminiproblems-problemphyxmini0270-recordedanswerchoice}
  \lean{PhyXMiniProblems.ProblemPhyXMini0270.recordedAnswerChoice}
  \uses{def:physics:phyx-mini-0270:phyxminiproblems-problemphyxmini0270-answerchoice}
  The answer label recorded by the supplied dataset
\end{definition}
\begin{proof}
  This is the declared model definition of recorded Answer Choice.
\end{proof}
\begin{definition}[Matches Displayed Torsion Constant]
  \label{def:physics:phyx-mini-0270:phyxminiproblems-problemphyxmini0270-matchesdisplayedtorsionconstant}
  \lean{PhyXMiniProblems.ProblemPhyXMini0270.MatchesDisplayedTorsionConstant}
  \uses{def:physics:phyx-mini-0270:phyxminiproblems-problemphyxmini0270-torsionconstantquantity, def:physics:phyx-mini-0270:phyxminiproblems-problemphyxmini0270-torsionconstantinnewtonmetersperradian, def:physics:phyx-mini-0270:phyxminiproblems-problemphyxmini0270-answerchoice}
  Agreement with a torsion constant displayed to the nearest tenth
\end{definition}
\begin{proof}
  This is the declared model definition of Matches Displayed Torsion Constant.
\end{proof}
\begin{definition}[Is Unique Closest Torsion Constant Choice]
  \label{def:physics:phyx-mini-0270:phyxminiproblems-problemphyxmini0270-isuniqueclosesttorsionconstantchoice}
  \lean{PhyXMiniProblems.ProblemPhyXMini0270.IsUniqueClosestTorsionConstantChoice}
  \uses{def:physics:phyx-mini-0270:phyxminiproblems-problemphyxmini0270-torsionconstantquantity, def:physics:phyx-mini-0270:phyxminiproblems-problemphyxmini0270-torsionconstantinnewtonmetersperradian, def:physics:phyx-mini-0270:phyxminiproblems-problemphyxmini0270-answerchoice}
  The selected value is strictly closer than every distinct displayed value
\end{definition}
\begin{proof}
  This is the declared model definition of Is Unique Closest Torsion Constant Choice.
\end{proof}
% --- Archon declaration topology end ---
% --- Archon physics formalization source end ---
